% \part{Aspectos Gerais}
\chapter[Introdução]{Introdução}

Estas instruções apresentam um conjunto mínimo de exigências necessárias a 
uniformidade de apresentação do relatório de Trabalho de Conclusão de Curso 
da FGA. Estilo, concisão e clareza ficam inteiramente sob a 
responsabilidade do(s) aluno(s) autor(es) do relatório.

As disciplinas de Trabalho de Conclusão de Curso (TCC) 01 e Trabalho de 
Conclusão de Curso (TCC) 02 se desenvolvem de acordo com Regulamento 
próprio aprovado pelo Colegiado da FGA. Os alunos matriculados nessas 
disciplinas devem estar plenamente cientes de tal Regulamento. 

\section{Motivação}

A

\section{Problema e Objetivo}

B

\section{Metodologia}

C

\subsection{Cronograma}

O planejamento das atividades a serem desenvolvidas ao longo do TCC 1 é
apresentado na Tabela \ref{tab:cronograma}, que organiza as etapas do
trabalho em intervalos de datas e associa a cada uma delas o respectivo
objetivo a ser alcançado.

\begin{table}[htb]
\caption{Cronograma de atividades do TCC 1}
\label{tab:cronograma}
\centering
\begin{tabularx}{\textwidth}{cc>{\raggedright\arraybackslash}X}
\hline
\textbf{Data Inicial} & \textbf{Data Final} & \textbf{Objetivo(s)} \\
\hline
10/08/2026 & 31/08/2026 & 1 - Conversa inicial com o orientador;\newline 2 - Definição do escopo do projeto. \\
\hline
01/09/2026 & 18/09/2026 & 1 - Entendimento do sistema existente no Ecossistema CEDIS;\newline 2 - Levantamento da arquitetura do sistema existente. \\
\hline
19/09/2026 & 30/09/2026 & 1 - Definição do objetivo da pesquisa;\newline 2 - Levantamento da base teórica. \\
\hline
01/10/2026 & 30/10/2026 & 1 - Levantamento de requisitos;\newline 2 - Diagramação da arquitetura do módulo de identificação do domínio de conceitos; \newline 3 - Priorização dos requisitos; \newline 4 - Definição do MVP. \\
\hline
01/11/2026 & 16/11/2026 & 1 - Desenvolvimento do protótipo de alta fidelidade; \newline 2 - Construção do cronograma de desenvolvimento em Sprints. \\
\hline
17/11/2026 & 13/12/2026 & 1- Revisão e ajustes dos artefatos \textit{(diagramas, cronogramas, etc)} produzidos; \newline 2 - Finalização da monografia do TCC 1. \\
\hline
\end{tabularx}
\legend{Fonte: Dos Autores.}
\end{table}

